\documentclass[10pt,a4paper]{amsart}
\usepackage[margin=19mm]{geometry}
\usepackage{amsmath,amssymb,mathtools}
\usepackage[scr=rsfs,cal=euler]{mathalfa}
\usepackage{xfrac}
\usepackage[unicode,hidelinks,bookmarks=false]{hyperref}
\newcommand{\ie}{i.e.\ }
\newcommand{\cf}{cf.\ }
\newcommand{\resp}{respectively\ }
\newcommand{\Subsection}{\subsection}
\providecommand{\nobreakdash}{\nobreak\hbox{-}}
\setlength{\emergencystretch}{2em}
\multlinegap0pt
\usepackage{bm}
\usepackage{bbm}
\usepackage{dsfont}
\usepackage{xspace}
\usepackage{cancel}
\usepackage{mathtools}
\usepackage{amsthm}
\makeatletter
\@ifundefined{theorem}{\newtheorem{theorem}{Theorem}}{}
\@ifundefined{lemma}{\newtheorem{lemma}[theorem]{Lemma}}{}
\@ifundefined{remark}{\newtheorem{remark}[theorem]{Remark}}{}
\makeatother
\title[Construction of Solutions with Extraordinary Gradient Amplif]{Construction of Solutions with Extraordinary Gradient Amplification and Localization for Schr\"odinger Equations}
\author{Huaian Diao, Xieling Fan, Hongyu Liu}
\date{Source: arXiv:2601.01389v2 (CC BY 4.0)}
\begin{document}
\maketitle

\noindent\emph{Source and licence.} Construction of Solutions with Extraordinary Gradient Amplification and Localization for Schr\"odinger Equations. Version arXiv:2601.01389v2 (CC BY 4.0), distributed under the Creative Commons Attribution 4.0 International licence, as recorded in the arXiv licence metadata for this exact version and shown on the abstract page https://arxiv.org/abs/2601.01389v2. Full rights notice: licenses/arxiv-2601.01389-CC-BY-4.0.txt.\par\smallskip

\noindent\textit{Editorial scope.} Counted unit: the Lemma whose source label is \texttt{Lemma 5.1} in the article (source0.tex lines 1347--1361, source label \texttt{Lemma 5.1}), delivered together with the article's own complete original proof of it (source0.tex lines 1363--1384, one \texttt{proof} environment of 22 lines). No mathematics is written, repaired or re-derived: statement and proof are the article's own text line for line. Required context retained uncounted: 1.1 (lines 269--275); 1.2 (lines 312--317); thm:main\_result (lines 715--752); cf 3.5 (lines 754--756); eq Hg (lines 927--929); important estimate for H\_g (lines 935--937); u\_0 (lines 939--941); nonlinear hyperbolic equation with source term and zero initial condition (lines 1312--1319); eq:F2 (lines 1339--1343). Every definition, display and label of those lines is retained unchanged. Omitted from this package and not redistributed: 1--268, 276--311, 318--714, 753--753, 757--926, 930--934, 938--938, 942--1311, 1320--1338, 1344--1346, 1362--1362, 1385--1783. Nothing inside a retained excerpt is omitted or altered: every retained body is delivered exactly as the article's lines are written, so no omission receipt of any kind accompanies this package. Citations: the retained excerpts carry no citation command at all, so no bibliography entry of the article is transcribed and no reference is redistributed. Reference adaptation: none was needed and none was made; every reference inside the delivered unit resolves inside it, so no unresolved reference or citation is produced. Printed slips kept literal: no printed word, symbol, hypothesis or number of the counted statement or of its proof was altered. Layout and class: the article's own class is \texttt{amsart} with packages including tikz, graphicx, subfigure, amssymb, epstopdf, geometry, bm, amsmath, amsfonts, amsthm, mathrsfs, cite, graphics, framed; the wrapper is the lane's pinned \texttt{amsart} editorial wrapper at 10pt on A4, which restates the theorem-like environments the retained text needs and adds the article's own macro definitions that the retained text needs (none), with extra wrapper packages (bm, bbm, dsfont, xspace, cancel, mathtools). The delivered numbering is the unit's own. Assets: no retained span contains a figure environment, a graphics-inclusion command, a PDF-page inclusion, a table or a TikZ picture; no image file is copied into this package, the package declares no asset dependency and no image file is redistributed with it. Licence: version arXiv:2601.01389v2 of this article is distributed under the Creative Commons Attribution 4.0 International licence, as recorded in the arXiv licence metadata for this exact version and shown on the abstract page https://arxiv.org/abs/2601.01389v2. A full rights notice is delivered at licenses/arxiv-2601.01389-CC-BY-4.0.txt. Characters: every retained excerpt is pure ASCII, so the delivered engine, XeLaTeX with its default Latin Modern text font, reports no missing character.
\begin{equation}\label{1.1}
    \begin{cases}
        \mathrm{i}\partial_t u + \mathcal{L}_1 u = 0 & \text{in } \Omega \times (0,T], \\
        u(\mathbf{x},0) = \phi(\mathbf{x}) & \text{in } \Omega, \\
        u(\mathbf{x},t) = \psi(\mathbf{x},t) & \text{on } \partial\Omega \times [0,T],
    \end{cases}
\end{equation}

\begin{equation}\label{1.2}
    \begin{cases}
        \mathrm{i}\partial_t U + \mathcal{L}_2 U + N(U) = 0 & 	\text{in } \mathbb{R}^{d} 	\times (0,T], \\
        U(\mathbf{x},0) = \Phi(\mathbf{x}) & 	\text{in } \mathbb{R}^d,
    \end{cases}
\end{equation}

	\begin{theorem}\label{thm:main_result}
		For any predetermined positive parameters $\varepsilon < 1$, $r_0$ sufficiently small, and $n$ distinct points $\mathbf{x}_1,\dots,\mathbf{x}_n \in \partial D$, we can properly choose corresponding $n$ balls satisfying:
		\begin{enumerate}
			\renewcommand{\labelenumi}{(\roman{enumi})}
			\item For problem \eqref{1.1} $($bounded domain $\Omega$ with $D \Subset \Omega):$ $B_{r_0}(\mathbf{y}_i) \Subset \Omega \setminus \overline{D}$,
			\item For problem \eqref{1.2} $($full space $\mathbb{R}^d ):$ $B_{r_0}(\mathbf{y}_i) \Subset \mathbb{R}^d \setminus \overline{D}$,
		\end{enumerate}
		such that $|\mathbf{x}_i - \mathbf{y}_i| = 2r_0$, the balls are mutually disjoint and
		$$
		\operatorname{dist}(\partial D, B_{r_0}(\mathbf{y}_i))>\frac{r_0}{4},\ \forall i=1,2...,n.
		$$
	Then, there exists an integer $M := M(r_0)$ and a smooth function $u_0$ with the following properties:
		\begin{enumerate}
			\item $u_0$ has the form$:$
			\begin{equation*}
				u_0(\mathbf{x},t) = H_g(\mathbf{x})\exp(-\mathrm{i} t),
			\end{equation*}
			where $ H_g(\mathbf{x})$ is a Herglotz function given by \eqref{eq Hg}. 
			
			\item $u_0$ satisfies the free Schr\"odinger equation$:$
			\[
			\mathrm{i}\partial_{t} u_0 + \Delta u_0 = 0 \quad \text{in } \mathbb{R}^d \times [0,\infty).
			\]
			
			\item It exhibits small $C^{1,\frac{1}{2}}$-norm in $\overline{D}:$
			\begin{equation}\label{u0 smalles}
				||u_0(\cdot,t)||_{C^{1,\frac{1}{2}}(\overline{D})} < C_1\varepsilon \quad \text{for all}~ t \in [0,\infty),
			\end{equation}
			where constant $C_1$ depends on $D$ and $d$.
			
			\item For any predetermined integer $m \geqslant M(r_0,\omega)$, there exist points $\mathbf{y}^*_i \in \partial B_{r_0}(\mathbf{y}_i)$  ($i=1,\dots,n$) with large amplitude $u_0(\mathbf{y}^*_i,t)$ such that
			\begin{align}
				|u_0(\mathbf{y}^*_i,t)| &> \frac{\sqrt{2m+2}\,r_0^{-1}}{2\sqrt{2\pi}} - C_2\varepsilon, & d &= 2, \label{2 d} \\
				|u_0(\mathbf{y}^*_i,t)| &> \frac{\sqrt{2m+3}\,r_0^{-3/2}}{16} - C_2\varepsilon, & d &= 3, \label{3 d}
			\end{align}
			for all $t \in [0,\infty)$, where constant $C_2$ depends on $r_0$, $D$, and $d$.
		\end{enumerate}
	\end{theorem}

	\begin{remark}\label{cf 3.5}
		From now and remaining of this paper, we denote by $\varepsilon$, $r_0$ the parameters, and by $u_0$ the function, specified in Theorem \ref{thm:main_result}.
	\end{remark}

		\begin{equation}\label{eq Hg}
			H_g(\mathbf{x}) = \int_{\mathbb{S}^{d-1}} g(\theta) \exp(\mathrm{i} \mathbf{x} \cdot \theta)  d\theta
		\end{equation}

		\begin{equation}\label{important estimate for H_g}
			\| H_g - v \|_{H^5(B_{r_0}(\mathbf{y}_i))} < \varepsilon \quad (i=1,\dots,n) \quad \text{and} \quad \| H_g \|_{H^5(D)} < \varepsilon.
		\end{equation}

		\begin{equation}\label{u_0}
			u_0(\mathbf{x},t) = H_g(\mathbf{x}) \exp(-\mathrm{i}  t).
		\end{equation}

	\begin{equation}\label{nonlinear hyperbolic equation with source term and zero initial condition}
		\left\{
		\begin{array}{@{}l@{\quad}l}
			\mathrm{i}\partial_t \mathscr{U}+\hat{\mathcal{L}}_2 \mathscr{U} + \hat{N}(\mathscr{U}) = F_2 & \text{in } \mathbb{R}^d \times (0,T], \\
			\mathscr{U}(\mathbf{x},0) = 0 & \text{in } \mathbb{R}^{d},
		\end{array}
		\right.
	\end{equation}

	\begin{equation}\label{eq:F2}
		F_2(\mathbf{x},t)
		= \nabla\cdot\big(\mathbf{I}-\mathscr{A}_2\big)\nabla u_0
		- c\,u_0 - \sum_{k=2}^{l_0}\alpha_k u_0^k.
	\end{equation}

	\begin{lemma}\label{Lemma 5.1}
		The coefficients, nonlinear term and source term in equation \eqref{nonlinear hyperbolic equation with source term and zero initial condition} satisfy the following properties:
		\begin{enumerate}
			\item $\operatorname{supp}\hat{c}(\cdot,t)$, $\operatorname{supp}\hat{N}(\mathscr{U})(\cdot,t)$
			and $\operatorname{supp}F_2(\cdot,t)\subset D$ for all $t\in[0,T]$.
			
			\item For a.e.\ $t\in[0,T]$,
			\[
			\|\hat{c}(\cdot,t)\|_{W^{2,\infty}(D)}\le C,\qquad
			\|F_2(\cdot,t)\|_{H^{3}(D)} \le C\,\varepsilon,
			\]
			where the constant $C$ depends only on $d$, $l_0$, $D$, $\|\mathscr{A}_2\|_{L^\infty(0,T)}$,
			$\|c\|_{L^\infty(0,T;W^{3,\infty}(D))}$, and $\|\alpha_k\|_{L^\infty(0,T;W^{3,\infty}(D))}$.
		\end{enumerate}
	\end{lemma}

	\begin{proof}
		Properties 1 follows immediately from the compact support in $D$ of $\mathbf{I} - \mathscr{A}_2(\cdot, t)$, $c(\cdot, t)$, $\alpha_k(\cdot, t)$ for $t \in [0, T]$, $k = 2, \dots, l_0$ combined with the expressions defining $\hat{c}$, $\hat{N}(\mathscr{U})$ and $F_2$.
		
		For properties 2, note that $\alpha_k$ are smooth, so all their derivatives are bounded in $D$. It remains to show $\partial_{\mathbf{x}}^{\alpha} u_0 \in L^{\infty}(D)$ for all multi-indices $\alpha$ with $|\alpha| \leqslant 3$. By the Sobolev embedding theorem, \eqref{important estimate for H_g}, and \eqref{u_0}, we have
		\begin{equation}\label{5.2}
			\|\partial_{\mathbf{x}}^{\alpha} u_0\|_{L^{\infty}(D)} \leqslant C \|\partial_{\mathbf{x}}^{\alpha} u_0\|_{H^{2}(D)} \leqslant C \|H_g\|_{H^5(D)} \leqslant C \varepsilon \leqslant C, \quad \forall |\alpha| \leqslant 3, \quad \forall t \in [0, T],
		\end{equation}
		where the final inequality holds since $\varepsilon \leqslant 1$ (See Theorem \ref{thm:main_result} and Remark \ref{cf 3.5}). Consequently,
		\begin{equation*}
			\|\hat{c}(\cdot, t)\|_{W^{2,\infty}(D)}\leqslant  C \left( \|c_2(\cdot, t)\|_{W^{2,\infty}(D)} + \sum_{k=2}^{l_0} \|u_0\|_{W^{2,\infty}(D)}^{k-1} \right) \leqslant C \left( 1 + \sum_{k=1}^{l_0-1} \varepsilon^{k} \right) \leqslant C,
		\end{equation*}
		for a.e. $t \in [0, T]$.
		
		For the source term, using the Banach algebra property of $H^3(D)$ and
	\eqref{important estimate for H_g}, \eqref{eq:F2}, and since $\varepsilon\le1$,
	we obtain
	\begin{equation}\label{eq:F2_est}
		\|F_2(\cdot,t)\|_{H^3(D)} \le C\sum_{k=1}^{l_0}\|H_g\|_{H^5(D)}^k
		\le C\sum_{k=1}^{l_0}\varepsilon^k \le C\varepsilon,
	\end{equation}
	for a.e.\ $t\in[0,T]$.
	\end{proof}

\end{document}
